\documentclass[9pt,a4paper,twocolumn]{ctexart}

% --- Packages ---
\usepackage[top=1.2cm, bottom=1.2cm, left=1.2cm, right=1.2cm, columnsep=0.8cm]{geometry}
\usepackage{hyperref}
\usepackage{graphicx}
\usepackage{float}
\usepackage{longtable}
\usepackage{tabularx}
\usepackage{booktabs}
\usepackage{enumitem}
\usepackage{xcolor}
\usepackage{fancyhdr}
\usepackage{listings}
\usepackage{fontspec}
\usepackage{titlesec}
\usepackage{caption}

% --- Code block styling ---
\definecolor{codebg}{RGB}{245,245,245}
\definecolor{codeframe}{RGB}{220,220,220}
\definecolor{codekw}{RGB}{0,0,192}
\definecolor{codecomment}{RGB}{0,128,0}
\definecolor{codestring}{RGB}{163,21,21}
\definecolor{codenum}{RGB}{128,0,128}
\definecolor{codepunct}{RGB}{90,90,90}
\definecolor{badgegreen}{RGB}{34,139,34}
\definecolor{badgeblue}{RGB}{30,144,255}
\definecolor{badgeblack}{RGB}{50,50,50}

% --- Difficulty badges (rounded corners via tikz, pre-rendered in saveboxes) ---
\usepackage{tikz}
\newsavebox{\badgechubox}
\savebox{\badgechubox}{\tikz[baseline]{\node[fill=badgegreen,text=white,rounded corners=2.5pt,inner xsep=4pt,inner ysep=1.5pt,font=\sffamily\footnotesize\bfseries]{初级};}}
\newsavebox{\badgezhongbox}
\savebox{\badgezhongbox}{\tikz[baseline]{\node[fill=badgeblue,text=white,rounded corners=2.5pt,inner xsep=4pt,inner ysep=1.5pt,font=\sffamily\footnotesize\bfseries]{中级};}}
\newsavebox{\badgegaobox}
\savebox{\badgegaobox}{\tikz[baseline]{\node[fill=badgeblack,text=white,rounded corners=2.5pt,inner xsep=4pt,inner ysep=1.5pt,font=\sffamily\footnotesize\bfseries]{高级};}}
\newcommand{\lvchu}{\usebox{\badgechubox}}
\newcommand{\lvzhong}{\usebox{\badgezhongbox}}
\newcommand{\lvgao}{\usebox{\badgegaobox}}
\pdfstringdefDisableCommands{%
  \def\lvchu{[初级]}%
  \def\lvzhong{[中级]}%
  \def\lvgao{[高级]}%
}

\lstset{
  basicstyle=\ttfamily\scriptsize,
  keywordstyle=\color{codekw}\bfseries,
  commentstyle=\color{codecomment}\itshape,
  stringstyle=\color{codestring},
  backgroundcolor=\color{codebg},
  frame=single,
  rulecolor=\color{codeframe},
  breaklines=true,
  breakatwhitespace=false,
  breakindent=0pt,
  postbreak=\mbox{\textcolor{red}{$\hookrightarrow$}\space},
  numbers=left,
  numberstyle=\tiny\color{gray},
  columns=flexible,
  keepspaces=true,
  showstringspaces=false,
  tabsize=2,
  aboveskip=4pt,
  belowskip=4pt,
  extendedchars=true,
  literate={，}{{，}}1 {；}{{；}}1 {！}{{！}}1 {？}{{？}}1 {“}{{“}}1 {”}{{”}}1 {（}{{（}}1 {）}{{）}}1 {《}{{《}}1 {》}{{》}}1,
}

% --- Extra language definitions (no built-in listings support) ---
\lstdefinelanguage{json}{
  morekeywords={true,false,null},
  sensitive=true,
  morestring=[b]",
  comment=[l]{//},
  morecomment=[s]{/*}{*/},
}
\lstdefinelanguage{yaml}{
  morekeywords={true,false,null,yes,no,on,off},
  sensitive=false,
  morecomment=[l]{\#},
  morestring=[b]",
  morestring=[d]',
}
\lstdefinelanguage{http}{
  morekeywords={GET,POST,PUT,DELETE,PATCH,HEAD,OPTIONS,HTTP,Host,Accept,Authorization,Content-Type,Content-Length,User-Agent},
  sensitive=true,
  morecomment=[l]{\#},
  morestring=[b]",
}
\lstdefinelanguage{TypeScript}{
  morekeywords={abstract,any,as,async,await,break,case,catch,class,const,continue,debugger,declare,default,delete,do,else,enum,export,extends,false,finally,for,from,function,get,if,implements,import,in,instanceof,interface,keyof,let,module,namespace,never,new,null,number,of,package,private,protected,public,readonly,return,set,static,string,super,switch,symbol,this,throw,true,try,type,typeof,undefined,unknown,var,void,while,with,yield,Record,Array,Promise,AsyncIterable,ReadableStream,Date,Error,Object,JSON},
  sensitive=true,
  morecomment=[l]{//},
  morecomment=[s]{/*}{*/},
  morestring=[b]",
  morestring=[b]',
  morestring=[b]`{`},
}

% --- Compact spacing ---
\linespread{0.95}
\setlength{\parskip}{1pt plus 1pt minus 1pt}
\setlength{\parindent}{0pt}

% --- Table styling ---
\renewcommand{\arraystretch}{1.1}

% --- Heading styling (numbered) ---
\titleformat{\section}{\normalsize\bfseries}{\thesection\quad}{0em}{}
\titlespacing{\section}{0pt}{6pt}{2pt}
\titleformat{\subsection}{\small\bfseries}{\thesubsection\quad}{0em}{}
\titlespacing{\subsection}{0pt}{4pt}{1pt}
\titleformat{\subsubsection}{\small\bfseries}{\thesubsubsection\quad}{0em}{}
\titlespacing{\subsubsection}{0pt}{3pt}{1pt}

% --- Page style ---
\pagestyle{fancy}
\fancyhf{}
\fancyhead[L]{\small\emph{\leftmark}}
\fancyhead[R]{\small\thepage}
\renewcommand{\headrulewidth}{0.4pt}
\renewcommand{\sectionmark}[1]{}  % prevent sections from overwriting leftmark

% --- Hyperref setup ---
\hypersetup{
  colorlinks=true,
  linkcolor=blue,
  urlcolor=blue,
  citecolor=blue,
}

% --- Caption format ---
\DeclareCaptionFormat{myformat}{\small\bfseries #1#2#3}
\captionsetup{format=myformat}

\begin{document}

\title{Agent 部署与发布：从灰度到回滚}
\date{}
\maketitle
\markboth{TypeScript 版 Agent 知识}{ TypeScript 版 Agent 知识}

\begin{quote}
语言策略：code（Java / Python / TypeScript）
\end{quote}



\section*{TL;DR}

\begin{enumerate}[itemsep=1pt, topsep=2pt]
  \item Agent 的发布单元不止代码：\textbf{Prompt、工具、模型路由、记忆策略都会改变线上行为}，必须一起版本化。
  \item 推荐路径：离线回归 → 影子流量 → 灰度 → 全量；每个阶段都有可观测指标和回滚条件。
  \item 部署前先回答三个问题：任务能不能幂等？跑一半断了怎么办？回滚时旧 Prompt 和旧工具是否还在？
\end{enumerate}


\section*{1. 部署形态}

{\footnotesize
\begin{tabular}{|p{\dimexpr 0.424\columnwidth-2\tabcolsep\relax}|p{\dimexpr 0.212\columnwidth-2\tabcolsep\relax}|p{\dimexpr 0.364\columnwidth-2\tabcolsep\relax}|}
\hline
\textbf{形态} & \textbf{适用场景} & \textbf{注意点} \\
\hline
同步 API 服务 & 交互式助手 & 需要流式、超时与取消 \\
\hline
异步任务队列 & 长任务、批处理 & 任务状态、幂等、死信队列 \\
\hline
定时/事件驱动 Worker & 巡检、告警处理 & 并发控制、去重 \\
\hline
边端/混合部署 & 数据不出域 & 模型版本与配置下发 \\
\hline
\end{tabular}
}



\section*{2. 发布单元与版本化}

{\footnotesize
\begin{tabular}{|p{\dimexpr 0.395\columnwidth-2\tabcolsep\relax}|p{\dimexpr 0.442\columnwidth-2\tabcolsep\relax}|p{\dimexpr 0.163\columnwidth-2\tabcolsep\relax}|}
\hline
\textbf{发布单元} & \textbf{版本标识} & \textbf{回滚方式} \\
\hline
应用代码 & Git commit / 镜像 tag & 回滚镜像 \\
\hline
System Prompt & prompt\_version & 配置回滚 \\
\hline
工具 Schema / Skill & toolset\_version & 注册表回滚 \\
\hline
模型路由策略 & route\_version & 网关配置回滚 \\
\hline
评测集与门禁 & eval\_version & 随发布记录关联 \\
\hline
\end{tabular}
}



发布包必须记录以上五个版本，Trace 中也要带上 \texttt{prompt\_version} 与 \texttt{toolset\_version}，否则线上问题无法归因。



\section*{3. 发布流水线}


\begin{lstlisting}
代码/配置变更
   │
   ▼
离线回归（成功率和工具准确率门禁）
   │
   ▼
影子流量（只记录，不影响用户）
   │
   ▼
灰度 5% → 25% → 100%
   │
   ▼
观察 P95、成功率、成本、危险动作率
   │
   ▼
通过则全量，否则自动回滚
\end{lstlisting}



灰度期间高危工具默认只读或审批升级；灰度时间不少于一个完整业务周期。



\section*{4. 运行时要求}

\begin{itemize}[itemsep=1pt, topsep=2pt]
  \item 健康检查：\texttt{/health} 同时检查模型供应商、向量库和工具依赖；
  \item 超时与取消：Agent Loop 总超时、单工具超时、用户取消；
  \item 幂等：同一 \texttt{task\_id} 重复提交只执行一次；
  \item 限流：按用户、租户、工具设置配额；
  \item 审计：部署事件、配置变更、回滚事件全部留痕。
\end{itemize}


\section*{5. TypeScript 版服务骨架与容器化}


\begin{lstlisting}[language=TypeScript]
interface AgentRunRequest {
  taskId: string;
  input: string;
  promptVersion: string;
  userId: string;
}

interface AgentRunResponse {
  taskId: string;
  status: "ok" | "timeout" | "error";
  result: string;
  elapsedMs: number;
}

interface HealthStatus {
  status: "up" | "degraded" | "down";
  modelUp: boolean;
  vectorDbUp: boolean;
  toolsUp: boolean;
}

interface AgentRuntime {
  run(request: AgentRunRequest): Promise<AgentRunResponse>;
  health(): HealthStatus;
}

import { Hono } from "hono";

export function buildApp(runtime: AgentRuntime) {
  const app = new Hono();

  app.post("/run", async (c) => {
    const body = await c.req.json<AgentRunRequest>();
    return c.json(await runtime.run(body));
  });

  app.get("/health", (c) => c.json(runtime.health()));

  return app;
}
\end{lstlisting}



\begin{lstlisting}
FROM node:22-alpine
WORKDIR /app
COPY package*.json ./
RUN npm ci --omit=dev
COPY . .
EXPOSE 3000
CMD ["node", "dist/index.js"]
\end{lstlisting}



\begin{lstlisting}
健康检查：GET /health
发布命令：docker build -t agent-service:<git-sha> . && docker push agent-service:<git-sha>
回滚命令：kubectl set image deploy/agent-service agent-service=agent-service:<previous-sha>
\end{lstlisting}



\section*{6. 回滚决策表}

{\footnotesize
\begin{tabular}{|p{\dimexpr 0.231\columnwidth-2\tabcolsep\relax}|p{\dimexpr 0.500\columnwidth-2\tabcolsep\relax}|p{\dimexpr 0.269\columnwidth-2\tabcolsep\relax}|}
\hline
\textbf{信号} & \textbf{阈值示例} & \textbf{动作} \\
\hline
任务成功率 & 低于基线 5 个百分点 & 立即回滚 \\
\hline
P95 延迟 & 超过 SLO 的 120\\% & 回滚或扩容 \\
\hline
危险动作率 & 出现一次未授权高危动作 & 冻结 + 回滚 \\
\hline
单任务成本 & 超过预算 130\\% & 回滚并复盘 \\
\hline
用户投诉 & 同主题短时激增 & 暂停灰度 \\
\hline
\end{tabular}
}



\section*{7. 上线 Runbook 检查单}

\begin{itemize}[itemsep=1pt, topsep=2pt]
  \item [ ] 发布包记录了代码、Prompt、工具、模型路由、评测集五个版本；
  \item [ ] 离线回归门禁通过；
  \item [ ] 回滚镜像和旧配置可用；
  \item [ ] 高危工具已配置审批；
  \item [ ] 监控面板和告警已就绪；
  \item [ ] 灰度比例、时长、回滚触发条件已明确；
  \item [ ] 事故联系人已确认。
\end{itemize}


\section*{8. 延伸阅读}

\begin{itemize}[itemsep=1pt, topsep=2pt]
  \item \href{04-大模型网关.md}{大模型网关}
  \item \href{05-Agent可观测与追踪.md}{Agent 可观测与追踪}
  \item \href{../07-系统设计/01-AI应用系统设计.md}{AI 应用系统设计}
\end{itemize}

\end{document}
